
\subsection{Verified Example: Memoized Function Calls} \label{sec:memoize}

\begin{figure}

\noindent\begin{minipage}{.5\textwidth}

\begin{lstlisting}[language=Verus,style=VerusLineNos]
#[spec] fn expected_result() -> u64 { /* ... */ }

#[exec] fn computation() -> u64 {
    ensures(|res: u64| res == expected_result()); %*\lclnum\label{clnum:comp:post}*
    /* ... */
}

#[spec] fn cell_value_inv(v: Option<u64>) -> bool {
    equal(v, Option::Some(expected_result()))
    || equal(v, Option::None) %*\lclnum\label{clnum:memo:cell-value-inv}*
}

#[spec] fn cell_is_valid(
        cell: InvCell<Option<u64>>) -> bool {
    cell.wf()
    && forall(|v| (#[trigger] cell.inv(v) ==
        cell_value_inv(v)))
}

#[exec] fn init_cell() -> InvCell<Option<u64>> {
    ensures(|c| cell_is_valid(c));
    InvCell::new(Option::None,
        |v: Option<u64>| cell_value_inv(v)) %*\lclnum\label{clnum:memo:init}*
}
\end{lstlisting}
\end{minipage}\hfill
\noindent\begin{minipage}{.5\textwidth}

\begin{lstlisting}[language=Verus,style=VerusLineNos]
#[exec] fn memoized_computation(
        cell: &InvCell<Option<u64>>) -> u64 {
    requires(cell_is_valid(*cell));
    ensures(|res: u64| res == expected_result());

    match cell.get() {
        Option::Some(res) => res, %*\lclnum\label{clnum:memo:case-none}*
        Option::None => {
            let res = computation();
            cell.replace(Option::Some(res));
            res %*\lclnum\label{clnum:memo:case-some}*
        }
    }
}

struct Client<'a> {
    cell: &'a InvCell<Option<u64>>,
}

fn main() {
    let c = init_cell();
    let client1 = Client { cell: &c };
    let client2 = Client { cell: &c };
    let x = memoized_computation(&client1.cell);
    let y = memoized_computation(&client2.cell);
    assert(x == y);
}
\end{lstlisting}
\end{minipage}

    \caption{
        \label{fig:memoize}
        Memoization example built on top of \texttt{InvCell}. %We take as given some (spec) value \texttt{expected\_result()} and
        %some (exec) code \texttt{computation()} that returns \texttt{expected\_result()}.
        %We then produce a function \texttt{memoized\_computation(cell)} which also returns
        %\texttt{expected\_result()}.
    }
\end{figure}

\emph{Memoization} is an optimization technique whereby a user saves time
by storing the result of a computation the first time it is invoked; on future invocations,
they use the stored value. Here, we show how to memoize a function call
\texttt{computation()}.
In \autoref{fig:memoize} we use a function that takes 0 arguments for simplicity, so there is only a single value to memoize;
in our supplementary materials~\cite{verussupplementary}, we provide a slightly more complex example that memoizes
a single-argument function \texttt{computation(i)}.

To set up the problem, we assume that \verb`computation()` has a postcondition~\lclref{clnum:comp:post} ensuring
that its result is equal to some desired (spec-mode) value, \verb`expected_result()`.
(This is similar to the setup of \verb`fibo_impl` and \verb`fibo` from earlier.)
The aim is to construct a function \verb`memoized_computation` that also returns
\verb`expected_result()`.
To keep the problem interesting, we also insist that it be possible to share the 
``result store'' across potentially many clients.
As such, we need to use a shared reference type \verb`&`; however,
a given update invocation might need to update the result store, which requires
mutability. Therefore, we need to use some form of interior mutability.

In our approach, we use
an \verb`InvCell` with a simple invariant on the data held by the cell.
When initializing the cell~\lclref{clnum:memo:init}, we specify the data invariant as a (spec-mode) boolean predicate on the interior values;
here, we set it to the function \verb`cell_value_inv`, defined at \lclref{clnum:memo:cell-value-inv}.
The resulting property of the cell is expressed as in \verb`cell_is_valid`.
This definition says that the value is valid if and only if the stored value is either
\verb`None` (not yet computed) or contains the correct answer.

To implement \verb`memoized_computation`, we first read from the cell; if the value we get is
\verb`Some`, then we return the value immediately~\lclref{clnum:memo:case-none} (as we can assume it satisfies the invariant
we just specified).
Otherwise, we perform the computation, store it in the cell,
and return it~\lclref{clnum:memo:case-some}.

Finally, in \verb`main`, we show that we can create multiple ``clients,'' sharing a
reference to the cell, and use them to call \verb`memoized_computation`.
